\section{Some Important Properties of Modulo and Division}
\label{sec:properties}

This chapter develops the concrete identities that follow from the
definitions and the linear-shift invariant of Section~5. It establishes:

\begin{itemize}
  \item The base cases where normalization is immediate: a small dividend,
    self-division, division by one, and agreement with the native modulo
    operator (Subsections~6.1--6.4)
  \item How a single or repeated shift of the dividend by the divisor moves
    the quotient without disturbing the remainder (Subsections~6.5--6.6)
  \item The uniqueness of the normalized remainder and its idempotence under
    repeated reduction (Subsections~6.7--6.8)
  \item Distributivity of modulo and division over addition and subtraction
    (Subsections~6.9--6.10)
  \item Shift invariance when the dividend is already divisible by the base,
    and the symmetry of remainder pairs around that base
    (Subsections~6.11--6.12)
  \item The unit-step increment law and the density of zero remainders across
    consecutive integers (Subsections~6.13--6.14)
\end{itemize}

\subsection{Trivial Case}
\label{sec:trivial-case}

If the dividend is smaller than a positive divisor, the candidate state
$\operatorname{DivMod}(a,b,0,a)$ is already final. No subtraction of $b$ is
needed, so the quotient is zero and the remainder is the original dividend.

\begin{equation*}
\begin{aligned}
& \forall\, a, b \in \mathbb{N} : b \neq 0 \\
& a < b \implies a \operatorname{mod} b = a \\
& a < b \implies a \operatorname{div} b = 0
\end{aligned}
\end{equation*}

This property is verified in
\href{https://github.com/thiagomata/prime-numbers/blob/master/src/main/scala/v1/chapter2/div/properties/ModSmallDividend.scala}{\texttt{ModSmallDividend}}.

\subsection{Identity}
\label{sec:identity}

The modulo of every number by itself is zero and the division of every number
by itself is one.

\begin{equation*}
\begin{aligned}
\forall\, n \in \mathbb{N} : n &\neq 0 \\
n \operatorname{mod} n &= 0 \\
n \operatorname{div} n &= 1
\end{aligned}
\end{equation*}

The candidate state normalizes in one shift to the final state with quotient
$1$ and remainder $0$:

\begin{equation*}
\begin{aligned}
n &= n\cdot 0+n = n\cdot 1+0 \\
0 &\leq 0 < \lvert n\rvert \\
\operatorname{DivMod}(n,n,0,n).\operatorname{solve}
  &= \operatorname{DivMod}(n,n,1,0). \quad \blacksquare\ \text{[Q.E.D.]}
\end{aligned}
\end{equation*}

This property is verified in
\href{https://github.com/thiagomata/prime-numbers/blob/master/src/main/scala/v1/chapter2/div/properties/ModIdentity.scala}{\texttt{ModIdentity::modIdentity}}.
A longer source proof showing the normalization path is available in
\href{https://github.com/thiagomata/prime-numbers/blob/master/src/main/scala/v1/chapter2/div/properties/ModIdentity.scala#longProof}{\texttt{ModIdentity::longProof}}.

\subsection{Modulo and Division by One}
\label{sec:mod-one}

Modulo by one always returns zero and division by one always returns the
dividend.

\begin{equation*}
\begin{aligned}
\forall\, n \in \mathbb{N} &: \\
n \operatorname{mod} 1 &= 0 \\
n \operatorname{div} 1 &= n
\end{aligned}
\end{equation*}

Both identities follow directly from the already canonical decomposition
$n=1\cdot n+0$; no induction or later unit-step result is needed.

\begin{equation*}
\begin{aligned}
n &= 1\cdot n+0,\qquad 0\leq 0<1 \\
\operatorname{mod}(n,1) &= 0,\qquad \operatorname{div}(n,1)=n.
  \quad \blacksquare\ \text{[Q.E.D.]}
\end{aligned}
\end{equation*}

These properties are verified in
\href{https://github.com/thiagomata/prime-numbers/blob/master/src/main/scala/v1/chapter2/div/properties/ModOne.scala}{\texttt{ModOne::modOneIsZero}}
and
\href{https://github.com/thiagomata/prime-numbers/blob/master/src/main/scala/v1/chapter2/div/properties/ModOne.scala}{\texttt{ModOne::divOneIsN}}.

\subsection{Compatibility with Native Modulo}
\label{sec:native-compatibility}

For non-negative dividends and positive divisors, the recursively normalized
modulo agrees with \texttt{BigInt}'s native \texttt{\%} operator.

\begin{equation*}
\begin{aligned}
\forall\, a, b \in \mathbb{Z} &: a \geq 0,\; b > 0 \\
a \operatorname{mod} b &= a \mathbin{\%} b
\end{aligned}
\end{equation*}

This is not a new mathematical fact but a bridge lemma: it confirms that the
native \texttt{\%} operator and the recursively defined
$\operatorname{mod}$ agree on their shared domain of non-negative dividends
and positive divisors, so results derived from one match results derived from
the other.

This property is verified in
\href{https://github.com/thiagomata/prime-numbers/blob/master/src/main/scala/v1/chapter2/div/properties/ModNativeCompatibility.scala#percentEqualsCalcMod}{\texttt{ModNativeCompatibility::percentEqualsCalcMod}}.

\subsection{Quotient Invariance Under Linear Shift}
\label{sec:linear-shift}

Adding or subtracting the divisor from the dividend changes the quotient by
one but leaves the remainder unchanged.

\begin{equation*}
\begin{aligned}
\forall a,b,q,r \in \mathbb{Z} &: b \neq 0,\; a = bq + r \\
\operatorname{mod}(a + b, b) &= \operatorname{mod}(a, b) \\
\operatorname{div}(a + b, b) &= \operatorname{div}(a, b) + 1 \\
\operatorname{mod}(a - b, b) &= \operatorname{mod}(a, b) \\
\operatorname{div}(a - b, b) &= \operatorname{div}(a, b) - 1
\end{aligned}
\end{equation*}

Let $q=\operatorname{div}(a,b)$ and $r=\operatorname{mod}(a,b)$. The
normalized state satisfies $a=bq+r$, with $r$ already in the canonical
remainder interval. Shifting the dividend by one divisor gives two equally
canonical states:

\begin{equation*}
\begin{aligned}
a+b &= bq+r+b = b(q+1)+r \\
a-b &= bq+r-b = b(q-1)+r \\
\operatorname{DivMod}(a+b,b,q+1,r).\operatorname{solve}
  &= \operatorname{DivMod}(a+b,b,q+1,r) \\
\operatorname{DivMod}(a-b,b,q-1,r).\operatorname{solve}
  &= \operatorname{DivMod}(a-b,b,q-1,r).
  \quad \blacksquare\ \text{[Q.E.D.]}
\end{aligned}
\end{equation*}

This property is verified for the
\href{https://github.com/thiagomata/prime-numbers/blob/master/src/main/scala/v1/chapter2/div/properties/AdditionAndMultiplication.scala#APlusBSameModPlusDiv}{positive case}
and
\href{https://github.com/thiagomata/prime-numbers/blob/master/src/main/scala/v1/chapter2/div/properties/AdditionAndMultiplication.scala#ALessBSameModDecreaseDiv}{negative case}.

\subsection{Quotient Invariance Under Linear Shift by Multiplier}
\label{sec:linear-shift-multiplier}

Adding a multiple of the divisor changes the quotient by that multiplier but
leaves the remainder unchanged.

As a direct consequence of the one-step shift laws, we can also prove that:

\begin{equation*}
\begin{aligned}
\forall a,b,q,r,m \in \mathbb{Z} &: b \neq 0,\; a = bq + r \\
\operatorname{mod}(a + m \cdot b, b) &= \operatorname{mod}(a, b) \\
\operatorname{div}(a + m \cdot b, b) &= \operatorname{div}(a, b) + m \\
\operatorname{mod}(a - m \cdot b, b) &= \operatorname{mod}(a, b) \\
\operatorname{div}(a - m \cdot b, b) &= \operatorname{div}(a, b) - m
\end{aligned}
\end{equation*}

For $m\geq0$, induction applies the positive one-step law to $a+m b$: the
remainder is unchanged at each step and the quotient gains one. Thus the
result holds at $m=0$ and is preserved from $m$ to $m+1$. The subtraction
identities follow by the same induction using the negative one-step law.
If $m<0$, write $m=-t$ with $t>0$ and exchange the addition and subtraction
identities just obtained.

\begin{equation*}
\begin{aligned}
\operatorname{mod}(a+(m+1)b,b)
  &= \operatorname{mod}((a+mb)+b,b) = \operatorname{mod}(a,b) \\
\operatorname{div}(a+(m+1)b,b)
  &= \operatorname{div}(a+mb,b)+1 = \operatorname{div}(a,b)+m+1.
  \quad \blacksquare\ \text{[Q.E.D.]}
\end{aligned}
\end{equation*}

This property is verified for the
\href{https://github.com/thiagomata/prime-numbers/blob/master/src/main/scala/v1/chapter2/div/properties/AdditionAndMultiplication.scala#APlusMultipleTimesBSameMod}{positive case}
and
\href{https://github.com/thiagomata/prime-numbers/blob/master/src/main/scala/v1/chapter2/div/properties/AdditionAndMultiplication.scala#ALessMultipleTimesBSameMod}{negative case}.

\subsection{Unique Remainder}
\label{sec:unique-remainder}

There is only one single remainder value for every $a, b$ pair with $b > 0$.

\begin{equation*}
\begin{aligned}
\forall\, a, b &\in \mathbb{N} : b > 0 \\
\exists!\, r \in \mathbb{N} &:
  0 \leq r < b \;\land\;
  a = \left\lfloor \frac{a}{b} \right\rfloor \cdot b + r
\end{aligned}
\end{equation*}

in other words, two $\operatorname{DivMod}$ instances with the same dividend
$a$ and divisor $b$ will have the same solution.

\begin{equation*}
\begin{aligned}
\forall a,b,q_x,r_x,q_y,r_y & \in \mathbb{N}, \\
\text{where } b & \neq 0 \text{, } \\
a & = bq_x + r_x \text{ and } \\
a & = bq_y + r_y \text{ then } \\
\operatorname{DivMod}(a,b,q_x,r_x).\operatorname{solve}
  & = \operatorname{DivMod}(a,b,q_y,r_y).\operatorname{solve}
\end{aligned}
\end{equation*}

For every $a,b$ pair, with any candidate quotients and remainders
$(q_x,r_x)$ and $(q_y,r_y)$ representing the same dividend, normalization
reaches the same solution. This property is verified in
\href{https://github.com/thiagomata/prime-numbers/blob/master/src/main/scala/v1/chapter2/div/properties/ModIdempotence.scala#modUnique}{\texttt{ModIdempotence::modUnique}}.

\subsection{Modulo Idempotence}
\label{sec:modulo-idempotence}

Taking the modulo of a number twice gives the same result as taking it once.

\begin{equation*}
\begin{aligned}
\forall\, a, b \in \mathbb{Z} &: b \neq 0 \\
a \operatorname{mod} b
  &= ( a \operatorname{mod} b ) \operatorname{mod} b
\end{aligned}
\end{equation*}

This property is verified in
\href{https://github.com/thiagomata/prime-numbers/blob/master/src/main/scala/v1/chapter2/div/properties/ModIdempotence.scala#modIdempotence}{\texttt{ModIdempotence::modIdempotence}}.

\subsection{Distributivity over Addition}
\label{sec:distributivity-addition}

The modulo operation distributes over addition, meaning that the remainder
of a sum equals the remainder of the sum of remainders. This allows us to
break down complex modulo operations into simpler components.

\begin{equation*}
\begin{aligned}
( a + c ) \operatorname{mod} b
  &= ( a \operatorname{mod} b + c \operatorname{mod} b )
     \operatorname{mod} b \\
( a + c ) \operatorname{div} b
  &= a \operatorname{div} b + c \operatorname{div} b
     + ( a \operatorname{mod} b + c \operatorname{mod} b )
       \operatorname{div} b \\
( a + c) \operatorname{mod} b
  &= (a \operatorname{mod} b) + (c \operatorname{mod} b)
     - b \cdot (((a \operatorname{mod} b)
     + (c \operatorname{mod} b)) \operatorname{div} b)
\end{aligned}
\end{equation*}

Let $q_a=\operatorname{div}(a,b)$, $r_a=\operatorname{mod}(a,b)$,
$q_c=\operatorname{div}(c,b)$, and $r_c=\operatorname{mod}(c,b)$. Thus
$a=bq_a+r_a$ and $c=bq_c+r_c$. Normalizing the sum of the two remainders
gives $r_a+r_c=bs+t$, where $s=\operatorname{div}(r_a+r_c,b)$ and
$t=\operatorname{mod}(r_a+r_c,b)$. Substitution gives the quotient and
remainder of $a+c$:

\begin{equation*}
\begin{aligned}
a+c &= b(q_a+q_c)+(r_a+r_c) \\
    &= b(q_a+q_c+s)+t \\
t &= r_a+r_c-b\cdot\operatorname{div}(r_a+r_c,b).
  \quad \blacksquare\ \text{[Q.E.D.]}
\end{aligned}
\end{equation*}

This property is verified in
\href{https://github.com/thiagomata/prime-numbers/blob/master/src/main/scala/v1/chapter2/div/properties/ModOperations.scala#modAdd}{\texttt{ModOperations::modAdd}}.
The third identity, isolating the multiple of $b$ subtracted out, is proved
directly in
\href{https://github.com/thiagomata/prime-numbers/blob/master/src/main/scala/v1/chapter2/div/properties/ModIdempotence.scala#modModPlus}{\texttt{ModIdempotence.scala\#modModPlus}}.

\subsection{Distribution over Subtraction}
\label{sec:distribution-subtraction}

Similar to addition, the modulo operation distributes over subtraction. The
remainder of a difference equals the remainder of the difference of
remainders, with appropriate handling of negative values.

\begin{equation*}
\begin{aligned}
( a - c ) \operatorname{mod} b
  &= ( a \operatorname{mod} b - c \operatorname{mod} b )
     \operatorname{mod} b \\
( a - c ) \operatorname{div} b
  &= a \operatorname{div} b - c \operatorname{div} b
     + ( a \operatorname{mod} b - c \operatorname{mod} b )
       \operatorname{div} b \\
( a - c ) \operatorname{mod} b
  &= (a \operatorname{mod} b) - (c \operatorname{mod} b)
     - b \cdot (((a \operatorname{mod} b)
     - (c \operatorname{mod} b)) \operatorname{div} b)
\end{aligned}
\end{equation*}

Let $q_a=\operatorname{div}(a,b)$, $r_a=\operatorname{mod}(a,b)$,
$q_c=\operatorname{div}(c,b)$, and $r_c=\operatorname{mod}(c,b)$. Thus
$a=bq_a+r_a$ and $c=bq_c+r_c$. Normalize the difference $r_a-r_c=bs+t$.
Then the canonical decomposition of $a-c$ is obtained without any assumption
that $r_a-r_c$ is nonnegative:

\begin{equation*}
\begin{aligned}
a-c &= b(q_a-q_c)+(r_a-r_c) \\
    &= b(q_a-q_c+s)+t \\
t &= r_a-r_c-b\cdot\operatorname{div}(r_a-r_c,b).
  \quad \blacksquare\ \text{[Q.E.D.]}
\end{aligned}
\end{equation*}

This property is verified in
\href{https://github.com/thiagomata/prime-numbers/blob/master/src/main/scala/v1/chapter2/div/properties/ModOperations.scala#modLess}{\texttt{ModOperations::modLess}}.
The third identity, isolating the multiple of $b$ subtracted out, is proved
directly in
\href{https://github.com/thiagomata/prime-numbers/blob/master/src/main/scala/v1/chapter2/div/properties/ModIdempotence.scala#modModMinus}{\texttt{ModIdempotence.scala\#modModMinus}}.

\subsection{Modular Shift Invariance under Divisible Base}
\label{sec:divisible-base-shift}

When a number is a multiple of the divisor (modulo equals zero), adding any
value does not change the modulo of that value. This property simplifies
calculations when one operand is already divisible by the base. It holds for
any integer $c$, including negative values.

\begin{equation*}
\begin{aligned}
\forall\, a, b, c \in \mathbb{Z} &: b \neq 0 \\
a \operatorname{mod} b = 0
  &\implies ( a + c ) \operatorname{mod} b = c \operatorname{mod} b
\end{aligned}
\end{equation*}

The addition law from Subsection~6.9 reduces the left-hand side to the
modulo of the two remainders. The hypothesis removes the first one, and
modulo idempotence removes the remaining repetition:

\begin{equation*}
\begin{aligned}
\operatorname{mod}(a+c,b)
  &= \operatorname{mod}(\operatorname{mod}(a,b)+\operatorname{mod}(c,b),b) \\
  &= \operatorname{mod}(\operatorname{mod}(c,b),b) \\
  &= \operatorname{mod}(c,b). \quad \blacksquare\ \text{[Q.E.D.]}
\end{aligned}
\end{equation*}

This property is verified in
\href{https://github.com/thiagomata/prime-numbers/blob/master/src/main/scala/v1/chapter2/div/properties/ModOperations.scala#modZeroPlusC}{\texttt{ModOperations::modZeroPlusC}}.

Substituting $-c$ for $c$ gives the subtraction corollary directly, since
$c$ is unrestricted:

\begin{equation*}
\begin{aligned}
\forall\, a, b, c \in \mathbb{Z} &: b \neq 0 \\
a \operatorname{mod} b = 0
  &\implies ( a - c ) \operatorname{mod} b = ( -c ) \operatorname{mod} b
\end{aligned}
\end{equation*}

This corollary is verified by the same lemma,
\href{https://github.com/thiagomata/prime-numbers/blob/master/src/main/scala/v1/chapter2/div/properties/ModOperations.scala#modZeroPlusC}{\texttt{ModOperations::modZeroPlusC}},
called with $-c$ in place of $c$.

\subsection{Symmetrical Modulo Pairs}
\label{sec:symmetrical-modulo-pairs}

The modulo of a value and the modulo of its complement relative to the base
sum to the base.

\begin{equation*}
\begin{aligned}
& b > 0,\quad 0 < k < b \\
& k \operatorname{mod} b + (b - k) \operatorname{mod} b = b
\end{aligned}
\end{equation*}

Since both $k$ and $b-k$ already lie inside the canonical remainder
interval, their remainders are themselves:

\begin{equation*}
\begin{aligned}
0<k<b &\implies \operatorname{mod}(k,b)=k \\
0<b-k<b &\implies \operatorname{mod}(b-k,b)=b-k \\
\operatorname{mod}(k,b)+\operatorname{mod}(b-k,b)
  &= k+(b-k)=b. \quad \blacksquare\ \text{[Q.E.D.]}
\end{aligned}
\end{equation*}

This property is verified in
\href{https://github.com/thiagomata/prime-numbers/blob/master/src/main/scala/v1/chapter2/div/properties/ModSum.scala}{\texttt{ModSum::sumSymmetricalMods}}.
The source excerpt is included in Appendix~A.2.

\subsection{Unit-Step Modulo-Division Increment Law}
\label{sec:unit-step}

When incrementing a number by one, the modulo cycles from 0 to b-1 and
resets, while the division increments only when the modulo reaches its
maximum value. This captures the "carry" behavior of division when counting.

\begin{equation*}
\begin{aligned}
\forall\, a, b \in \mathbb{N} : b \neq 0 \\
a \operatorname{mod} b = b - 1
  &\implies (a + 1) \operatorname{mod} b = 0 \\
a \operatorname{mod} b \neq b - 1
  &\implies (a + 1) \operatorname{mod} b = (a \operatorname{mod} b) + 1 \\
a \operatorname{mod} b = b - 1
  &\implies (a + 1) \operatorname{div} b = (a \operatorname{div} b) + 1 \\
a \operatorname{mod} b \neq b - 1
  &\implies (a + 1) \operatorname{div} b = a \operatorname{div} b
\end{aligned}
\end{equation*}

Let $q=\operatorname{div}(a,b)$ and $r=\operatorname{mod}(a,b)$, so
$a=bq+r$ with $0\leq r<b$. If $r=b-1$, adding one produces the canonical
state $(q+1,0)$. Otherwise $r<b-1$, so $(q,r+1)$ is already canonical. This
also covers $b=1$: only the first case can occur.

\begin{equation*}
\begin{aligned}
r=b-1 &\implies a+1=bq+(b-1)+1=b(q+1)+0 \\
r<b-1 &\implies a+1=bq+(r+1),\quad 0\leq r+1<b.
  \quad \blacksquare\ \text{[Q.E.D.]}
\end{aligned}
\end{equation*}

This property is verified in
\href{https://github.com/thiagomata/prime-numbers/blob/master/src/main/scala/v1/chapter2/div/properties/ModOperations.scala#addOne}{\texttt{ModOperations::addOne}}.

\subsection{Consecutive Integers: Zero Density}
\label{sec:zero-density}

In any block of $p$ consecutive integers, exactly one is divisible by $p$.
This is the basic counting form of modulo periodicity: as we advance through
consecutive integers, the remainder modulo $p$ visits zero once per complete
period.

\textbf{At most one zero per block.} If $\operatorname{mod}(a,p)=0$ and
$0 < d < p$, then $\operatorname{mod}(a+d,p)\neq 0$. Within any block of
size $p$ starting from a multiple, no later offset inside the same block can
also be divisible by $p$.

\begin{equation*}
\begin{aligned}
&\operatorname{mod}(a,\; p) = 0 \;\land\; 0 < d < p
  \implies \operatorname{mod}(a + d,\; p) \neq 0
  \quad \text{[nonzeroAfterZero]}
\end{aligned}
\end{equation*}

\textbf{At least one zero per block.} For any starting value $n$ and modulus
$p>1$, there exists a $k \in [0,p)$ such that
$\operatorname{mod}(n+k,p)=0$. The witness is $k=0$ when $n$ is already
divisible by $p$, and $k=p-\operatorname{mod}(n,p)$ otherwise.

\begin{equation*}
\begin{aligned}
&\forall n \geq 0,\; p > 1,\; \exists\, k \in [0, p) :
  \operatorname{mod}(n + k,\; p) = 0
  \quad \text{[existsZero]}
\end{aligned}
\end{equation*}

\textbf{Exactly one zero per block.} Existence gives a zero offset, while
uniqueness says two zero offsets in the same block must be equal. Together,
among $p$ consecutive integers starting from $n$, exactly one is divisible
by $p$.

\begin{equation*}
\begin{aligned}
&\forall n \geq 0,\; p > 1,\; \exists!\, k \in [0, p) :
  \operatorname{mod}(n + k,\; p) = 0
  \quad \text{[existsZero + atMostOneZero]}
\end{aligned}
\end{equation*}

More explicitly, let $k$ be the offset supplied by existence. If $i$ and
$j$ are two offsets in $[0,p)$ with zero remainder, the at-most-one result
applied to $i$ and $j$ yields $i=j$; hence the existing $k$ is unique.

\begin{equation*}
\begin{aligned}
&\exists\, k\in[0,p):\ \operatorname{mod}(n+k,p)=0 \\
&\operatorname{mod}(n+i,p)=\operatorname{mod}(n+j,p)=0,\quad i,j\in[0,p)
  \implies i=j \\
&\therefore\ \exists!\, k\in[0,p):\ \operatorname{mod}(n+k,p)=0.
  \quad \blacksquare\ \text{[Q.E.D.]}
\end{aligned}
\end{equation*}

{\raggedright
These properties are verified in
\href{https://github.com/thiagomata/prime-numbers/blob/master/src/main/scala/v1/chapter2/div/properties/ConsecutiveIntegers.scala}{\texttt{ConsecutiveIntegers::\allowbreak nonzeroAfterZero}},
\href{https://github.com/thiagomata/prime-numbers/blob/master/src/main/scala/v1/chapter2/div/properties/ConsecutiveIntegers.scala}{\texttt{ConsecutiveIntegers::\allowbreak existsZero}},
\href{https://github.com/thiagomata/prime-numbers/blob/master/src/main/scala/v1/chapter2/div/properties/ConsecutiveIntegers.scala}{\texttt{ConsecutiveIntegers::\allowbreak exactlyOneZeroInConsecutive}},
and
\href{https://github.com/thiagomata/prime-numbers/blob/master/src/main/scala/v1/chapter2/div/properties/ConsecutiveIntegers.scala}{\texttt{ConsecutiveIntegers::\allowbreak atMostOneZero}}.
The compact source shape is included in Appendix~A.3. The maintained source
is
\href{https://github.com/thiagomata/prime-numbers/blob/master/src/main/scala/v1/chapter2/div/properties/ConsecutiveIntegers.scala}{\texttt{ConsecutiveIntegers.scala}}.
The file also contains multi-factor density helpers, such as
\texttt{twoFactorsDensity} and \texttt{densityForFactorList}; their
statements are not established in this article, and the multiplicative
density extension they aim at is discussed as an open direction in the
Future Work section.
\par}




